\lecture

\begin{guiding}
For which base spaces does the slogan ``bundles are the same as maps to the
Grassmannian'' hold, and what does the Grassmannian look like as a space? The
technical hypothesis is paracompactness; the geometric answer is that the
Grassmannian is a CW complex with exactly one cell for each Schubert symbol.
\end{guiding}

\section{Paracompactness}

\begin{definition}
A space is \emph{paracompact} if every open cover admits a locally finite
refinement, that is, a refinement such that every point has a neighbourhood
meeting only finitely many of its members.
\end{definition}

\begin{theorem}[quoted]
Every metric space is paracompact and Hausdorff. A regular space which is a
countable union of compact subsets is paracompact and Hausdorff. Consequently
the direct limit of a sequence of compact Hausdorff spaces is paracompact.
\end{theorem}

\begin{remark}
The infinite Grassmannian $G_n=\varinjlim_k G_n(\R^{n+k})$ is therefore
paracompact, being a direct limit of compact spaces. All bases we shall meet ---
manifolds, CW complexes, Grassmannians --- are paracompact.
\end{remark}

The following technical lemma replaces an arbitrary trivializing cover by a
countable one. It is the device that makes the classification theorem work for
non-compact bases.

\begin{lemma}\label{lem:countable-cover}
Let $\xi$ be a vector bundle over a paracompact Hausdorff base $B$. Then there
is a countable open cover $U_1,U_2,\dots$ of $B$ such that $\xi|_{U_k}$ is
trivial for every $k$.
\end{lemma}

\begin{proof}
\begin{strategy}
Start from a locally finite trivializing cover $\{V_\alpha\}$ with subordinate
bump functions $\lambda_\alpha$. For each \emph{finite} index set $S$ collect the
points where the functions indexed by $S$ dominate all the others; this produces
an open set contained in $\bigcap_{\alpha\in S}V_\alpha$, hence trivializing.
Two different sets of the same cardinality give disjoint pieces, so we may take
their union and index the resulting cover by the cardinality $k=|S|$ alone.
\end{strategy}

\step{1} \emph{Set-up.} Choose a locally finite open cover $\{V_\alpha\}$ of $B$
such that $\xi|_{V_\alpha}$ is trivial, and continuous functions
$\lambda_\alpha\colon B\to[0,1]$ with $\lambda_\alpha=0$ outside $V_\alpha$ and
with the open sets $W_\alpha=\lambda_\alpha^{-1}\big((0,1]\big)$ still covering
$B$. Such data exists by paracompactness.

\step{2} \emph{The sets $U(S)$.} For a finite non-empty set $S$ of indices put
\[
U(S)=\Big\{b\in B\ \Big|\ \min_{\alpha\in S}\lambda_\alpha(b)
\ >\ \max_{\beta\notin S}\lambda_\beta(b)\Big\}.
\]
Near any point only finitely many $\lambda$'s are non-zero, so both the minimum
and the maximum are continuous there, and $U(S)$ is open.

\step{3} \emph{$U(S)$ is trivializing.} If $b\in U(S)$ and $\alpha\in S$ then
$\lambda_\alpha(b)>0$, so $b\in V_\alpha$. Hence
\[
U(S)\ \subset\ \bigcap_{\alpha\in S}V_\alpha\ \subset\ V_\alpha ,
\]
and $\xi$ restricted to $U(S)$ is trivial.

\step{4} \emph{Sets of equal size are disjoint.} Let $S\neq S'$ with
$|S|=|S'|=k$. Choose $\alpha\in S\smallsetminus S'$ and
$\alpha'\in S'\smallsetminus S$. If $b$ were in $U(S)\cap U(S')$ we would get
$\lambda_\alpha(b)>\lambda_{\alpha'}(b)$ from $b\in U(S)$, since
$\alpha\in S$ and $\alpha'\notin S$, and symmetrically
$\lambda_{\alpha'}(b)>\lambda_\alpha(b)$ from $b\in U(S')$: a contradiction.

\step{5} \emph{The cover.} Put
\[
U_k=\bigsqcup_{|S|=k}U(S),
\]
a disjoint union of open trivializing sets by Steps 3 and 4, hence itself open
and trivializing, since a bundle trivial on each member of a disjoint family is
trivial on their union. Finally the $U_k$ cover $B$: given $b$, let
$S=\{\alpha\mid\lambda_\alpha(b)\ \text{is maximal}\}$, which is finite and
non-empty by local finiteness and Step 1; then $b\in U(S)\subset U_{|S|}$.
\end{proof}

\section{The universal bundle theorem}

\begin{theorem}\label{thm:universal}
Let $\xi$ be an $\R^n$-bundle over a paracompact base $B$. Then

\begin{enumerate}
\item there exists a bundle map $f\colon\xi\to\gamma^n$;
\item any two bundle maps $\xi\to\gamma^n$ are homotopic through bundle maps.
\end{enumerate}

Consequently the homotopy class of the base map
$\bar f\in\Maps(B,G_n)$ depends only on $\xi$, and $\xi\iso\bar f^{\,*}\gamma^n$.
\end{theorem}

\begin{proof}
\begin{strategy}
Existence repeats the construction of Lemma~\ref{lem:compact-classify}, with the
countable cover of Lemma~\ref{lem:countable-cover} replacing the finite one; the
target is $\R^\infty$ instead of a fixed $\R^m$. For uniqueness, the naive
straight-line homotopy between two fibrewise injective maps
$\hat f,\hat g\colon E(\xi)\to\R^\infty$ works whenever their values are never
negative multiples of each other. In general one first pushes $\hat f$ into the
odd coordinates and $\hat g$ into the even ones, where the obstruction cannot
occur, and chains three such homotopies.
\end{strategy}

\step{1} \emph{Existence.} Let $U_1,U_2,\dots$ be as in
Lemma~\ref{lem:countable-cover}, with trivializations
$h_i\colon U_i\times\R^n\to\pi^{-1}(U_i)$ and bump functions $\lambda_i$ whose
positivity sets cover $B$. Setting
$\tilde h_i=\proj_{\R^n}\circ h_i^{-1}$ and
$h'_i(e)=\lambda_i(\pi(e))\,\tilde h_i(e)$, extended by zero, define
\[
\hat f\colon E(\xi)\longrightarrow \R^n\oplus\R^n\oplus\cdots=\R^\infty,
\qquad
\hat f(e)=\big(h'_1(e),\,h'_2(e),\,\dots\big).
\]
Locally only finitely many terms are non-zero, so $\hat f$ is continuous; it is
linear on fibres, and injective on the fibre over $b$ because
$\lambda_i(b)>0$ for some $i$. As in Step 1 of
Lemma~\ref{lem:compact-classify}, $\hat f$ determines a bundle map
$f\colon\xi\to\gamma^n$ by
$f(e)=\big(\hat f(F_{\pi(e)}(\xi)),\ \hat f(e)\big)$.

\step{2} \emph{The easy homotopy.} Suppose $\hat f,\hat g\colon E(\xi)\to
\R^\infty$ are continuous, linear and injective on fibres, and satisfy
\[
\hat f(e)\ \neq\ -\lambda^2\,\hat g(e)
\qquad\text{for all }e\neq0\text{ and all }\lambda\in\R .
\tag{$\ast$}
\]
Then $\hat h_t=(1-t)\hat f+t\hat g$ is fibrewise linear, and injective on
fibres: if $\hat h_t(v)=0$ with $v\neq0$ and $0<t<1$ then
$\hat f(v)=-\tfrac{t}{1-t}\hat g(v)$, contradicting $(\ast)$; for $t=0,1$ it is
$\hat f$ or $\hat g$. So $\hat h_t$ gives a homotopy of bundle maps from $f$ to
$g$.

\step{3} \emph{Two shift operators.} Let $d_1,d_2\colon\R^\infty\to\R^\infty$ be
the linear injections
\[
d_1(e_i)=e_{2i-1},
\qquad
d_2(e_i)=e_{2i}.
\]
Being injective and linear, each induces a bundle map $\gamma^n\to\gamma^n$
sending $(X,v)\mapsto(d_jX,d_jv)$; and $d_j\circ\hat f$ is again continuous,
linear and fibrewise injective.

\step{4} \emph{$\hat f$ and $d_1\hat f$ satisfy $(\ast)$.} Let $v=\hat f(e)\neq0$
and let $k$ be the index of its first non-zero coordinate. Suppose
$d_1v=-\lambda^2 v$ with $\lambda\neq0$. If $k=1$ then $(d_1v)_1=v_1$, so
$v_1=-\lambda^2v_1$ forces $v_1=0$, a contradiction. If $k>1$ then the first
non-zero coordinate of $d_1v$ sits in position $2k-1>k$, so $(d_1v)_k=0$, giving
$-\lambda^2v_k=0$ and again $v_k=0$. The same argument applies to $\hat g$ and
$d_2\hat g$, using $2k>k$.

\step{5} \emph{$d_1\hat f$ and $d_2\hat g$ satisfy $(\ast)$.} The vector
$d_1\hat f(e)$ is supported in odd coordinates and $d_2\hat g(e)$ in even ones.
If $d_1\hat f(e)=-\lambda^2d_2\hat g(e)$ then both sides vanish, so
$\hat f(e)=0$ and hence $e=0$.

\step{6} \emph{Chaining.} By Steps 2, 4 and 5,
\[
f\ \simeq\ d_1\circ f\ \simeq\ d_2\circ g\ \simeq\ g
\]
through bundle maps, which proves (2).

\step{7} \emph{Conclusion.} A bundle map $f\colon\xi\to\gamma^n$ gives
$\xi\iso\bar f^{\,*}\gamma^n$ by Lemma~\ref{lem:bundle-map-pullback}, and by (2)
the class $[\bar f]$ is well defined.
\end{proof}

\begin{remark}
The converse also holds: homotopic maps $B\to G_n$ pull $\gamma^n$ back to
isomorphic bundles, so
\[
\big\{\R^n\text{-bundles over }B\big\}\big/\!\iso
\ \longleftrightarrow\
[\,B,\,G_n\,].
\]
This gives a clean definition of characteristic classes. For
$c\in H^i(G_n)$ and a bundle $\xi$ with classifying map $f_\xi$, put
\[
c(\xi):=f_\xi^*\,c\ \in\ H^i(B).
\]
Naturality is automatic: if $\eta=g^*\xi$ then $f_\eta\simeq f_\xi\circ g$, so
$c(\eta)=g^*c(\xi)$.
\[
\begin{tikzcd}[row sep=large, column sep=large]
\eta \arrow[d] \arrow[r] & \xi \arrow[d] \arrow[r] & \gamma^n\arrow[d]\\
B_1 \arrow[r,"g"] \arrow[rr, bend right=20, "f_\eta"'] & B_2 \arrow[r,"f_\xi"] & G_n
\end{tikzcd}
\]
The whole theory therefore reduces to one computation: what is $H^*(G_n)$?
\end{remark}

\section{A cell structure on the Grassmannian}

\begin{definition}
A \emph{CW complex} is a Hausdorff space $K$ partitioned into disjoint open
cells $\{e_\alpha\}$ such that

\begin{enumerate}
\item each $e_\alpha$ carries a characteristic map
$f\colon D^{n(\alpha)}\to K$ taking the open disc homeomorphically onto
$e_\alpha$;
\item $\bar e_\alpha\smallsetminus e_\alpha$ is contained in a finite union of
cells of strictly lower dimension;
\item every point lies in a finite subcomplex;
\item $K$ carries the direct limit topology of its finite subcomplexes.
\end{enumerate}
\end{definition}

Fix the standard flag $\R^0\subset\R^1\subset\dots\subset\R^m$, where $\R^k$ is
spanned by the first $k$ coordinates, and let $X\subset\R^m$ be an $n$-plane.
Consider the sequence of integers
\[
0=\dim(X\cap\R^0)\ \le\ \dim(X\cap\R^1)\ \le\ \dots\ \le\ \dim(X\cap\R^m)=n .
\]

\begin{claim}\label{claim:jumps}
Consecutive terms differ by at most $1$.
\end{claim}

\begin{proof}
The sequence
\[
0\longrightarrow X\cap\R^{k-1}\longrightarrow X\cap\R^{k}
\ \xrightarrow{\ x\,\mapsto\, x_k\ }\ \R
\]
is exact: the kernel of the $k$-th coordinate function on $X\cap\R^k$ is exactly
$X\cap\R^{k-1}$. Hence
\[
\dim(X\cap\R^k)=\dim(X\cap\R^{k-1})+\rk(x\mapsto x_k)
\le\dim(X\cap\R^{k-1})+1 .
\]
\end{proof}

So the sequence increases by exactly $1$ at $n$ places, the \emph{jumps}
\[
1\le\sigma_1<\sigma_2<\dots<\sigma_n\le m,
\]
called the \emph{Schubert symbol} of $X$. Define
\[
e(\sigma)=\big\{X\in G_n(\R^m)\ \big|\
\dim(X\cap\R^{\sigma_i})=i \ \text{ and }\
\dim(X\cap\R^{\sigma_i-1})=i-1,
\]
\[
\text{for } i=1,\dots,n\big\}.
\]
Every $X$ lies in exactly one $e(\sigma)$, and there are $\binom{m}{n}$ symbols.

\begin{claim}\label{claim:half-space}
Let $H^{k}\subset\R^k\subset\R^m$ be the half-space
$\{x\in\R^k\mid x_k>0\}$. Then $X\in e(\sigma)$ if and only if $X$ has a basis
$x_1,\dots,x_n$ with $x_i\in H^{\sigma_i}$; moreover such a basis can be chosen
\emph{orthonormal}, and then it is unique.
\end{claim}

\begin{proof}
\begin{strategy}
One direction is a rank computation: a vector whose last non-zero coordinate is
in position $\sigma_i$ forces the dimension to jump there. The other direction
constructs the basis one vector at a time, taking at each jump the unique unit
vector orthogonal to the previous ones and normalized to have positive
$\sigma_i$-th coordinate.
\end{strategy}

\step{1} \emph{Sufficiency.} Suppose $X$ has a basis with $x_i\in H^{\sigma_i}$.
Then $x_1,\dots,x_i\in\R^{\sigma_i}$, so $\dim(X\cap\R^{\sigma_i})\ge i$; and
the $\sigma_i$-th coordinate of $x_i$ is non-zero while
$x_1,\dots,x_{i-1}\in\R^{\sigma_{i-1}}\subset\R^{\sigma_i-1}$, so the map
$X\cap\R^{\sigma_i}\to\R$, $x\mapsto x_{\sigma_i}$, has rank $1$ and
$\dim(X\cap\R^{\sigma_i})>\dim(X\cap\R^{\sigma_i-1})$. Counting the $n$ jumps
this forces $\dim(X\cap\R^{\sigma_i})=i$ exactly, so $X\in e(\sigma)$.

\step{2} \emph{Necessity, first vector.} Let $X\in e(\sigma)$. Then
$X\cap\R^{\sigma_1}$ is a line and $X\cap\R^{\sigma_1-1}=0$, so every non-zero
vector of that line has non-zero $\sigma_1$-th coordinate. There is exactly one
unit vector $x_1$ on it with positive $\sigma_1$-th coordinate, and
$x_1\in H^{\sigma_1}$.

\step{3} \emph{Necessity, induction.} Suppose orthonormal
$x_1,\dots,x_{i-1}$ have been found with $x_j\in H^{\sigma_j}$, spanning
$X\cap\R^{\sigma_{i-1}}$. The space $X\cap\R^{\sigma_i}$ has dimension $i$ and
contains $X\cap\R^{\sigma_{i-1}}$ as a hyperplane, so its orthogonal complement
inside $X\cap\R^{\sigma_i}$ is a line. Since
$\dim(X\cap\R^{\sigma_i-1})=i-1=\dim(X\cap\R^{\sigma_{i-1}})$, no non-zero
vector of that line lies in $\R^{\sigma_i-1}$; hence its two unit vectors have
non-zero $\sigma_i$-th coordinate, and exactly one of them, call it $x_i$, has
it positive. Then $x_i\in H^{\sigma_i}$, and $x_1,\dots,x_i$ is orthonormal.

\step{4} \emph{Uniqueness.} At each stage the choice was forced: the line was
determined by $X$ and the previous vectors, and the sign by the positivity
condition.
\end{proof}

In matrix terms: writing a basis of $X$ as the rows of an $n\times m$ matrix and
row reducing, $X\in e(\sigma)$ if and only if $X$ is the row space of a unique
matrix in reduced echelon form whose pivots occupy the columns
$\sigma_1,\dots,\sigma_n$,
\[
\begin{pmatrix}
*&\cdots&*&1&0&\cdots& & &0\\
*&\cdots&*&*&*&1&0&\cdots&0\\
\vdots& & & & & & &\ddots&\\
*&\cdots&*&*&*&*&\cdots&1&0\cdots 0
\end{pmatrix}
\]
The number of free entries is
\[
d(\sigma)=(\sigma_1-1)+(\sigma_2-2)+\dots+(\sigma_n-n),
\]
which will be the dimension of the cell. Following
Claim~\ref{claim:half-space}, set
\[
e'(\sigma)=V^0_n(\R^m)\cap\big(H^{\sigma_1}\times\dots\times H^{\sigma_n}\big),
\qquad
\bar e'(\sigma)=V^0_n(\R^m)\cap\big(\bar H^{\sigma_1}\times\dots\times
\bar H^{\sigma_n}\big),
\]
so that $q=\spann$ maps $e'(\sigma)$ bijectively onto $e(\sigma)$. In the next
lecture we prove that $\bar e'(\sigma)$ is a closed cell of dimension
$d(\sigma)$ and that these cells assemble into a CW structure.
